% ------------------------------------------------------------------
% FILE: sections/11_next_steps.tex
% Practical next steps for implementation and manuscript transition.
% \input{} into main.tex — do NOT wrap in \documentclass or \begin{document}.
% ------------------------------------------------------------------

\section{Next Steps}\label{sec:nextsteps}

This final section translates the methodological framework established
in the preceding sections into a concrete implementation plan.  It
specifies the software stack, hardware profiling protocol, compression
strategy, reproducibility safeguards, and a phased timeline leading to
manuscript preparation.


% ------------------------------------------------------------------
\subsection{Implementation Stack}
\label{subsec:impl_stack}
% ------------------------------------------------------------------

All experimental code will be developed in \textbf{Python~3.10+},
using the following library ecosystem:

\paragraph{Classical models.}
Logistic Regression, Linear~SVM, RBF~SVM, Random Forest, Gradient
Boosted Trees (via XGBoost), K-Means, and a baseline MLP will be
implemented using
\textbf{scikit-learn}~\cite{sklearn_docs}.  Scikit-learn provides a
uniform \texttt{fit}/\texttt{predict}/\texttt{score} API, built-in
cross-validation utilities, and deterministic seeding via
\texttt{random\_state}, which simplifies reproducibility.

\paragraph{Neural models.}
The 1D~CNN (\Cref{subsec:cnn1d}), ResNet-1D (\Cref{subsec:resnet1d}),
and Kolmogorov--Arnold Network (\Cref{subsec:kan}) will be implemented
in \textbf{PyTorch}, which offers fine-grained control over custom
architectures, training loops, and gradient manipulation.  The MLP
(\Cref{subsec:mlp}) may also be implemented in PyTorch if the
scikit-learn version proves insufficient for architectural
experimentation (e.g., custom regularisation schedules).

\paragraph{Keras clarification.}
\textbf{Keras is a high-level API for defining and training neural
network architectures.}  It is not an algorithm and therefore does not
appear in the ten-model comparison (\Cref{sec:algorithms}).  Keras
(now integrated into TensorFlow as \texttt{tf.keras} and also
available as a standalone multi-backend library) can be used to
implement any of the neural architectures discussed in this report ---
MLP, 1D~CNN, or ResNet-1D --- but the choice of Keras vs.\ raw
PyTorch is an engineering decision, not a methodological one.  If
Keras is used for any model, this will be noted in the implementation
log, and the resulting model must still satisfy all evaluation criteria
from \Cref{sec:eval} identically to its PyTorch counterpart.

\paragraph{Supporting libraries.}
\begin{itemize}
  \item \textbf{NumPy / SciPy}: array operations, signal processing
        (unwrapping, derivatives).
  \item \textbf{Pandas}: metadata management and result aggregation.
  \item \textbf{Matplotlib / Seaborn}: all visualisations (spectra,
        confusion matrices, degradation curves, calibration diagrams).
  \item \textbf{SHAP}: post-hoc feature attribution for tree ensembles
        and neural models.
  \item \textbf{Optuna} or \textbf{scikit-optimize}: Bayesian
        hyperparameter optimisation with nested cross-validation.
\end{itemize}


% ------------------------------------------------------------------
\subsection{Edge Deployment Frameworks}
\label{subsec:edge_frameworks}
% ------------------------------------------------------------------

Trained models must be converted to edge-compatible inference formats.
Three frameworks are considered:

\begin{enumerate}
  \item \textbf{TensorFlow Lite}~\cite{tensorflow_lite}: the primary
        target for TensorFlow/Keras-trained models.  Supports INT8
        post-training quantisation, delegate-based hardware
        acceleration (GPU, DSP, NPU), and deployment on ARM Cortex-M
        and Cortex-A targets via the TFLite Micro interpreter.
  \item \textbf{ONNX Runtime}~\cite{onnx_runtime}: the primary target
        for PyTorch-trained models.  ONNX (Open Neural Network
        Exchange) provides a framework-agnostic intermediate
        representation.  ONNX Runtime supports quantisation,
        graph optimisation, and execution providers for ARM, x86,
        and RISC-V targets.
  \item \textbf{PyTorch Mobile}: an alternative for PyTorch models
        that avoids the ONNX conversion step.  Supports selective
        operator loading to reduce binary size, and \texttt{torch.jit}
        scripting for ahead-of-time compilation.
\end{enumerate}

For scikit-learn tree ensembles, the inference engine is simpler: the
trained model can be exported as a set of if--else rules (e.g., via
\texttt{sklearn-porter} or manual code generation) and compiled
directly into C or C++ for bare-metal execution.  No deep-learning
runtime is required.


% ------------------------------------------------------------------
\subsection{The Edge Inference Pipeline}
\label{subsec:inference_pipeline}
% ------------------------------------------------------------------

While previous sections focus on model training and selection, it is critical to define how the final chosen model will actually execute inference in a production edge environment. The inference pipeline consists of four sequential steps executed directly on the microcontroller or edge device:

\begin{enumerate}
  \item \textbf{Data Acquisition and Parsing:} The edge device receives a raw frequency sweep from the connected sensor hardware. It extracts the $4001$ magnitude (and optionally phase) values, discarding malformed readings exactly as defined in the validity preprocessing stage (\Cref{sec:preproc}).
  \item \textbf{Frozen Normalisation:} The device applies the exact normalisation constants (e.g., $\mu_{\text{train}}$ and $\sigma_{\text{train}}$) that were computed during the training phase (\Cref{sec:norm}). These constants are hardcoded into the edge firmware to prevent data leakage or drift.
  \item \textbf{Forward Pass Execution:} The normalised vector is passed into the compressed model. Because the model has been converted to an edge-compatible format (e.g., INT8 via TensorFlow Lite Micro), the mathematical operations---whether they are dot products for Logistic Regression, tree traversals for Random Forests, or convolutions for a 1D CNN---are executed using highly efficient integer arithmetic. This respects the strict SRAM and Flash limits of the hardware.
  \item \textbf{Decision and Abstention:} The final output layer produces a confidence score for the three classes (\texttt{dry}, \texttt{wet}, \texttt{ice}). If the highest score exceeds a predefined threshold, a classification is made. If the model is uncertain, it triggers an abstention protocol rather than guessing blindly (\Cref{sec:eval}).
\end{enumerate}

\paragraph{Why this makes the edge model a viable option.}
The core reason the shortlisted models are viable for this task is that their inference requirements (the ``forward pass'') can be decoupled from their training requirements. While training requires significant compute and memory, the inference pipeline outlined above is extremely lightweight. By processing the $S_{11}$ data locally without cloud dependency, the system achieves real-time latency, preserves bandwidth, and maintains offline reliability.

% ------------------------------------------------------------------
\subsection{Hardware Profiling Protocol}
\label{subsec:hw_profiling}
% ------------------------------------------------------------------

Edge deployment claims must be validated on representative hardware.
The profiling protocol specifies:

\begin{enumerate}
  \item \textbf{Desktop baseline:} all models profiled on an x86-64
        CPU (single-threaded, no GPU) to establish a normalised
        comparison.
  \item \textbf{ARM Cortex-A target:} models profiled on a
        Cortex-A class board (e.g., Raspberry Pi~4, Coral Dev Board)
        to represent single-board-computer deployments.
  \item \textbf{ARM Cortex-M target (if available):} a subset of
        the smallest models profiled on a Cortex-M class
        microcontroller (e.g., STM32H7, nRF5340) to test
        bare-metal feasibility.
  \item \textbf{RISC-V target (optional):} if a RISC-V development
        board is available, profiling one or two models validates
        architecture portability.
\end{enumerate}

For each hardware target, the following measurements will be recorded:
%
\begin{itemize}
  \item Inference latency: mean and standard deviation over
        $\geq$\,\num{100} runs.
  \item Peak RAM usage during inference.
  \item Flash/storage footprint of the compiled model binary.
  \item Power consumption per inference (if instrumentation is
        available).
\end{itemize}


% ------------------------------------------------------------------
\subsection{Quantisation Strategy}
\label{subsec:quantisation}
% ------------------------------------------------------------------

Model compression via quantisation follows a two-stage protocol:

\paragraph{Stage 1: INT8 post-training quantisation (PTQ).}
The trained float32 model is quantised to INT8 using a representative
calibration dataset (a subset of the training data, disjoint from
the test set).  Macro~$F_1$ is re-evaluated on the test set.  If the
$F_1$ drop is $\leq 0.02$ (absolute), PTQ is accepted.

\paragraph{Stage 2: Quantisation-aware training (QAT).}
If PTQ causes an unacceptable $F_1$ drop ($> 0.02$), the model is
retrained with quantisation-aware training, which simulates
quantisation noise during forward passes.  QAT typically recovers
most of the accuracy lost by PTQ at the cost of additional training
time.

For tree-based models, quantisation is less relevant because
inference involves integer comparisons against thresholds rather
than floating-point arithmetic.  The primary compression lever for
trees is ensemble pruning (removing low-contribution trees or
limiting depth).


% ------------------------------------------------------------------
\subsection{Manuscript Transition Plan}
\label{subsec:manuscript_plan}
% ------------------------------------------------------------------

This report establishes the methodology; the manuscript will report
the results.  The transition follows this structure:

\begin{enumerate}
  \item \textbf{Introduction:} motivate the RF/microwave sensing
        application, state the classification problem, and summarise
        the methodological framework.
  \item \textbf{Methods:} compress the content of
        \Cref{sec:dataset,sec:repr,sec:preproc,sec:norm,sec:overfit,%
        sec:augment,sec:algorithms,sec:eval} into a concise methods
        section, citing this report as a supplementary technical
        reference where detail is needed.
  \item \textbf{Results:} present the filled evaluation tables
        (\Cref{tab:macro_f1_results,tab:perclass_recall,%
        tab:model_complexity}), confusion matrices, degradation
        curves, and ablation findings.
  \item \textbf{Discussion:} interpret results, compare the
        shortlisted models, discuss failure cases, and address
        limitations (small sample size, single sensor geometry,
        laboratory-only data).
  \item \textbf{Conclusion:} recommend the primary model for
        deployment and identify directions for future work
        (multi-sensor fusion, online adaptation, expanded
        environmental conditions).
\end{enumerate}


% ------------------------------------------------------------------
\subsection{Data Versioning and Reproducibility}
\label{subsec:reproducibility}
% ------------------------------------------------------------------

\paragraph{Dataset snapshot.}
Before any training begins, the dataset will be frozen and versioned:
%
\begin{itemize}
  \item Compute SHA-256 checksums for every \texttt{.s1p} file.
  \item Record the total file count per class
        ($601 + 501 + 501 = 1603$).
  \item Store the checksum manifest alongside the code repository.
  \item Any additions, removals, or modifications to the dataset
        after the snapshot constitute a new dataset version and
        require re-running the full pipeline.
\end{itemize}

\paragraph{Reproducibility safeguards.}
\begin{itemize}
  \item \textbf{Random seeds:} all stochastic operations (train--test
        splits, weight initialisation, dropout masks, data shuffling)
        will use a fixed global seed (e.g., \texttt{seed=42}) and
        per-component seeds where supported.
  \item \textbf{Deterministic operations:} PyTorch deterministic mode
        (\texttt{torch.use\_deterministic\_algorithms(True)}) will be
        enabled where feasible.  Known non-deterministic operations
        (e.g., certain CUDA atomics) will be documented.
  \item \textbf{Hyperparameter logging:} every training run will log
        the full hyperparameter configuration, library versions
        (via \texttt{pip freeze}), hardware identifier, and wall-clock
        training time.
  \item \textbf{Result archiving:} trained model weights, evaluation
        metrics, and generated figures will be archived with
        timestamped identifiers.
\end{itemize}


% ------------------------------------------------------------------
\subsection{Timeline Estimate}
\label{subsec:timeline}
% ------------------------------------------------------------------

The experimental phase is organised into five sequential stages:

\begin{enumerate}
  \item \textbf{Data validation} (1~week): parse all \num{1603} files,
        verify frequency-grid consistency, compute checksums, perform
        exploratory visualisation, and run K-Means structure probe.
  \item \textbf{Baseline models} (1--2~weeks): train and evaluate the
        five classical models (Logistic Regression, Linear~SVM,
        RBF~SVM, Random Forest, Gradient Boosted Trees) across all
        Route~A encodings.
  \item \textbf{Neural models} (2--3~weeks): implement, train, and
        evaluate MLP, 1D~CNN, ResNet-1D, and KAN.  Includes
        architecture search, regularisation tuning, and augmentation
        experiments.
  \item \textbf{Edge profiling} (1~week): convert shortlisted models
        to TFLite/ONNX, apply quantisation, and benchmark on target
        hardware.
  \item \textbf{Manuscript preparation} (2--3~weeks): compile results,
        generate publication-quality figures, write and revise the
        manuscript.
\end{enumerate}

\noindent
\textbf{Estimated total:} 7--10~weeks from the start of the
experimental phase.  Stages may overlap where dependencies permit
(e.g., manuscript writing can begin during edge profiling).


% ------------------------------------------------------------------
\subsection{Closing Statement}
\label{subsec:closing}
% ------------------------------------------------------------------

\begin{quote}
\itshape
This report defines the methodological contract that all subsequent
experimental work will follow.  Every preprocessing step, normalisation
choice, splitting protocol, evaluation metric, and decision gate
described herein constitutes a binding commitment.  Deviations from
this pipeline must be documented, justified, and disclosed in the
experimental log and the final manuscript.
\end{quote}
